%% USPSC-pre-textual-ICMC.tex
%% Dados do trabalho: tipo de documento, area de concentracao, preambulo, titulacao.
%%
%% ATENCAO: este arquivo e carregado pela classe USPSC ANTES do \usepackage[utf8]{inputenc}.
%% Por isso os acentos sao escritos com sequencias de escape do LaTeX (\^e, \'a, \c{c}, \~a)
%% e NAO em UTF-8 direto. Os demais arquivos do projeto usam UTF-8 normalmente.

\instituicao{Instituto de Ci\^encias Matem\'aticas e de Computa\c{c}\~ao, Universidade de S\~ao Paulo}
\unidade{INSTITUTO DE CI\^ENCIAS MATEM\'ATICAS E DE COMPUTA\c{C}\~AO}
\unidademin{Instituto de Ci\^encias Matem\'aticas e de Computa\c{c}\~ao}
\universidademin{Universidade de S\~ao Paulo}
\setorpos{SERVI\c{C}O DE P\'OS-GRADUA\c{C}\~AO DO ICMC-USP}

\notafolharosto{Vers\~ao original}
\notafolharostoadic{Original version}
%Para a versao revisada, comente as duas linhas acima e descomente as duas abaixo
%\notafolharosto{Vers\~ao revisada}
%\notafolharostoadic{Final version}

% ---
% dados complementares para CAPA e FOLHA DE ROSTO
% ---
\universidade{UNIVERSIDADE DE S\~AO PAULO}

\titulo{An\'alise Preditiva de Churn de Clientes utilizando Machine Learning}
\titleabstract{Predictive Analysis of Customer Churn using Machine Learning}
\tituloadic{Predictive Analysis of Customer Churn using Machine Learning}
\tituloresumo{An\'alise Preditiva de Churn de Clientes utilizando Machine Learning}

\autor{Fabiano Rodrigues Ugulino}
\autorficha{Ugulino, Fabiano}
\autorabr{UGULINO, F.}

% O codigo Cutter definitivo e atribuido pela Biblioteca Prof. Achille Bassi.
% Enquanto nao houver o definitivo, o campo fica em branco (linha ativa abaixo).
\cutter{ }
%\cutter{U26a}

\local{S\~ao Carlos}
\data{2026}

\renewcommand{\orientadorname}{Orientadora:}
\orientador{Profa. Dra. Cibele M. Russo}
\orientadoradic{Advisor: Profa. Dra. Cibele M. Russo}
\orientadorcorpoficha{orientadora Cibele M. Russo}
\orientadorficha{Russo, Cibele M., orient}

\notaautorizacao{AUTORIZO A REPRODU\c{C}\~AO E DIVULGA\c{C}\~AO TOTAL OU PARCIAL DESTE TRABALHO, POR QUALQUER MEIO CONVENCIONAL OU ELETR\^ONICO PARA FINS DE ESTUDO E PESQUISA, DESDE QUE CITADA A FONTE.}

\notabib{Ficha catalogr\'afica elaborada pela Biblioteca Prof. Achille Bassi, ICMC/USP, com os dados fornecidos pelo(a) autor(a)}

\newcommand{\programa}[1]{
% MBAIAp ==========================================================================
\ifthenelse{\equal{#1}{MBAIAp}}{
	\tipotrabalho{Monografia (MBA em Intelig\^encia Artificial e Big Data)}
	\tipotrabalhoabs{Monograph (MBA in Artificial Intelligence and Big Data)}
	\area{Intelig\^encia Artificial}
	\areaadic{Concentration area: Artificial Intelligence}
	\preambulo{Monografia apresentada ao Departamento de Ci\^encias de Computa\c{c}\~ao do Instituto de Ci\^encias Matem\'aticas e de Computa\c{c}\~ao, Universidade de S\~ao Paulo - ICMC/USP, como parte dos requisitos para obten\c{c}\~ao do t\'itulo de Especialista em Intelig\^encia Artificial e Big Data.}
	\preambuloadic{Monograph presented to the Departamento de Ci\^encias de Computa\c{c}\~ao do Instituto de Ci\^encias Matem\'aticas e de Computa\c{c}\~ao, Universidade de S\~ao Paulo - ICMC/USP, as part of the requirements for obtaining the title of Specialist in Artificial Intelligence and Big Data.}
	\notaficha{Monografia (MBA em Intelig\^encia Artificial e Big Data)}
	\notacapaicmc{Monografia - MBA em Intelig\^encia Artificial e Big Data}
    }{
% MBAIAe ==========================================================================
\ifthenelse{\equal{#1}{MBAIAe}}{
	\tipotrabalho{Monografia (MBA em Intelig\^encia Artificial e Big Data)}
	\tipotrabalhoabs{Monograph (MBA in Artificial Intelligence and Big Data)}
	\renewcommand{\areaname}{Concentration area:}
	\area{Artificial Intelligence and Big Data}
	\areaadic{\'Area de concentra\c{c}\~ao: Intelig\^encia Artificial e Big Data}
	\preambulo{Monograph presented to the Departamento de Ci\^encias de Computa\c{c}\~ao do Instituto de Ci\^encias Matem\'aticas e de Computa\c{c}\~ao, Universidade de S\~ao Paulo - ICMC/USP, as part of the requirements for obtaining the title of Specialist in Artificial Intelligence and Big Data.}
	\preambuloadic{Monografia apresentada ao Departamento de Ci\^encias de Computa\c{c}\~ao do Instituto de Ci\^encias Matem\'aticas e de Computa\c{c}\~ao, Universidade de S\~ao Paulo - ICMC/USP, como parte dos requisitos para obten\c{c}\~ao do t\'itulo de Especialista em Intelig\^encia Artificial e Big Data.}
	\notaficha{Monograph (MBA in Artificial Intelligence and Big Data)}
	\notacapaicmc{Monograph - MBA in Artificial Intelligence and Big Data}
    }{
% Outros ==========================================================================
	\tipotrabalho{Monografia (MBA)}
	\tipotrabalhoabs{Monograph (MBA)}
	\area{Nome da \'Area}
	\areaadic{Additional area}
	\preambulo{Monografia apresentada ao Instituto de Ci\^encias Matem\'aticas e de Computa\c{c}\~ao, Universidade de S\~ao Paulo - ICMC/USP, como parte dos requisitos para obten\c{c}\~ao do t\'itulo de Especialista.}
	\preambuloadic{Monograph presented to the Instituto de Ci\^encias Matem\'aticas e de Computa\c{c}\~ao, Universidade de S\~ao Paulo - ICMC/USP, as part of the requirements for obtaining the title of Specialist.}
	\notaficha{Monografia (MBA)}
	\notacapaicmc{Monografia - MBA}
        }}
}
